\chapter{What Structure Buys, and Where It Stops Paying}
\label{ch:em-results}

Chapter~\ref{ch:em-method} built an estimator that is exact in its inference
step and approximate only in its numerical representation and its optimisation
budget. Chapter~\ref{ch:em-kernel} measured what it costs to run. This chapter
asks what it is \emph{worth}, and the question is comparative: worth relative to
a network trained on the same clean trajectories and given none of the
structure. What each arm is allowed to observe of those trajectories differs,
and Section~\ref{sec:em-estimand} states how before any number is quoted.

\begin{quote}
How much does supplying the chain reduce denoising error at a fixed number of
clean trajectories; how much of that reduction survives against a baseline that is given
locality but not the model; how much mixture capacity does the estimator need;
and how far can the Markov assumption itself be wrong before the structure
stops paying?
\end{quote}

\noindent The four are answered, at different strengths, and
Section~\ref{sec:em-summary} states each at the strength its evidence supports.

\section{What is compared, and how it is scored}\label{sec:em-estimand}

\subsection{The estimand}

Every headline number in this chapter is a ratio of two relative $L^2$ score
errors.

Fix a diffusion time $t$ and a test set $\{x^{(j)}\}_{j=1}^{J}$ of noised
sequences drawn from the model of Section~\ref{sec:em-model}. Let
$\score^{\star}(x,t) = \nabla_x \log p_t(x)$ be the exact joint score, supplied
by the validated grid reference of Section~\ref{sec:l-grid} under the true
kernel. For an estimator $\widehat{\score}$ write
\begin{equation}
  \mathcal{E}\big[\widehat{\score}; t\big]
  \;=\;
  \Bigg(
  \frac{\sum_{j=1}^{J}\norm{\widehat{\score}(x^{(j)},t) - \score^{\star}(x^{(j)},t)}^2}
       {\sum_{j=1}^{J}\norm{\score^{\star}(x^{(j)},t)}^2}
  \Bigg)^{1/2},
  \label{eq:em-relerr}
\end{equation}
the ratio of two Euclidean norms over the whole test set, so that $0.1$ means
the estimate is wrong by a tenth of the signal's own size. It is scale free:
multiplying the data by a constant leaves it unchanged. Pooling over the noise
schedule $\mathcal{T}$ means pooling the two sums, not averaging the ratios,
\begin{equation}
  \mathcal{E}^2\big[\widehat{\score}\big]
  \;=\;
  \sum_{t \in \mathcal{T}} w_t \,\mathcal{E}^2\big[\widehat{\score}; t\big],
  \qquad
  w_t \;=\;
  \frac{\E\norm{\score^{\star}(X_t,t)}^2}{\sum_{s \in \mathcal{T}}\E\norm{\score^{\star}(X_s,s)}^2}.
  \label{eq:em-pooling}
\end{equation}
The weights are not a free choice and they are not uniform. In the
implementation they are the measured reference score energies of the frozen
test bundle, level by level; the variance-matched Gaussian chain explains
their shape in closed form. There, $\score^{\star} = -\Sigt^{-1}x$ with
$\Sigt = e^{-2t}\Sigz + \Deltat I$, so
\begin{equation}
  \E\norm{\score^{\star}(X_t,t)}^2 \;=\; \operatorname{tr}\Sigt^{-1},
  \qquad
  \Sigt^{-1} \to I \ \ (t \to \infty),
  \label{eq:em-weights}
\end{equation}
and at the low-noise end the limit is \emph{finite}: since
$\Sigz \succ 0$, $\Sigt \to \Sigz$ and
\begin{equation}
  \operatorname{tr}\Sigt^{-1}
  \;=\; \operatorname{tr}\Qz
  \;+\; 2t\,\big[\operatorname{tr}\Qz - \operatorname{tr}\Qz^2\big]
  \;+\; O(t^2)
  \qquad (t \downarrow 0),
  \label{eq:em-weights-smallt}
\end{equation}
with no $1/t$ singularity --- the quantity that does diverge like $1/(2t)$ is
the likelihood precision $e^{-2t}/\Deltat$ in latent coordinates, which is a
different object. The trace nevertheless falls steeply across the schedule
(from $\operatorname{tr}\Qz = 193.4$ at $t \to 0$ to $\nsites = 32$ at
$t \to \infty$ for this $\alpha$ and $\nsites$; the linear coefficient in
\eqref{eq:em-weights-smallt} is $-3140$), so the finite schedule still
weights its lower-noise levels most heavily.
Table~\ref{tab:em-locality} makes that concrete: under the closed-form
weights the four smallest of the twelve levels carry $53.6\%$ of the pooled
weight, and under the measured Laplace energies actually used for pooling the
same four carry $54.5\%$. This matters twice
below --- once when the window head's optimal radius is predicted, and once when
its residual error is discussed.

The two arms are the EM--BP estimator of Chapter~\ref{ch:em-method}, which is
given the Markov factorisation, and a denoising-score-matching network trained
by \eqref{eq:sm-dsm} with no structural information at all. Both are fitted
from the same $\nseq$ clean trajectories, and that --- not ``the same data''
--- is the equality the protocol enforces. What each arm is allowed to see of
those trajectories is not the same, and the difference runs \emph{against} the
result reported here, so it is stated plainly rather than left to the
appendix:

\begin{itemize}
  \item EM--BP sees one noised realisation $x$ of each trajectory and never
  sees the clean trajectory $a$ at all; the noised bundle is frozen before
  fitting begins.
  \item The network is trained on \emph{paired} $(a, x)$ data, as
  \eqref{eq:sm-dsm} requires, and redraws the diffusion time and the
  corrupting noise at every gradient step, so over a run it consumes many
  distinct noisy views of each clean trajectory rather than one.
\end{itemize}

\noindent
Supplying the Markov factorisation is thus traded against supplying the clean
signal and unlimited corruptions of it, and both asymmetries are deliberate:
the baseline is the standard training recipe for this model class, and
weakening it to match EM--BP's view would measure a handicap rather than a
structural advantage. The consequence for reading
\eqref{eq:em-ratio} is that the ratio is a statement about
\emph{value per clean trajectory} under the two standard recipes, not about
generic data efficiency at a matched information interface.
Table~\ref{tab:em-architectures} records what each arm in the wider ladder is
given.
Writing $\mathcal{E}_{\mathrm{base}}$ and $\mathcal{E}_{\mathrm{EM}}$ for
their errors on a common test set, the reported quantity is
\begin{equation}
  R \;=\; \frac{\mathcal{E}_{\mathrm{base}}}{\mathcal{E}_{\mathrm{EM}}},
  \qquad R > 1 \ \text{means the structured estimator is more accurate,}
  \label{eq:em-ratio}
\end{equation}
evaluated cell by cell, a cell being one (training-set size, seed) pair, with
both arms re-fitted independently in every cell.

One aggregation choice remains, and it matters because \eqref{eq:em-ratio} is a
ratio of the quantities \eqref{eq:em-relerr} and not a linear functional of
them. This chapter reports exactly the pooled definition above, paired at the
seed level: for each seed and size, squared error and reference energy are
summed over the twelve levels before the square root, giving one pooled error
per arm; the ratio \eqref{eq:em-ratio} is formed within the seed from those
two pooled errors; and the reported number is the mean of the sixteen paired
per-seed ratios, with its standard error. The seed is the inferential unit
throughout --- the twelve levels share a training set and a fitted model.
Appendix~\ref{app:aggregation} recomputes the same cells under four
alternative summaries; the direction and the order of magnitude are the same
in all of them, but the values are not interchangeable: the mean of
per-level ratios, which earlier drafts reported, is systematically the
largest ($\ratiocellwiselo$--$\ratiocellwisehi\times$ on these cells,
against the pooled $\ratiolo$--$\ratiohi\times$), because it lets levels
with little reference energy count as much as levels with a great deal.

\subsection{The common protocol}

Every comparison in Sections~\ref{sec:em-pointwise} to~\ref{sec:em-capacity}
runs on the same data-generating model and the same schedule.
Table~\ref{tab:em-protocol} fixes it once so that the sections need not repeat
it; Section~\ref{sec:em-nonmarkov}, which deliberately changes the prior, and
Section~\ref{sec:em-shape-information}, which is a likelihood calculation rather
than a fit, state their own departures.

\begin{center}
\begin{minipage}{\linewidth}\small\centering
\captionof{table}[The protocol shared by every comparison]{The protocol shared
by the comparisons of Sections~\ref{sec:em-pointwise}--\ref{sec:em-capacity}.
The generating prior is a Laplace AR(1) chain, so the reference score is
genuinely nonlinear (Chapter~\ref{ch:laplace}) and no arm is being compared
against a linear target. Validation, test and training bundles are disjoint,
and no arm sees the test bundle before its budget is chosen.}
\label{tab:em-protocol}
\begin{tabular}{@{}ll@{}}
\toprule
generating prior & Laplace AR(1), $\corr = 0.85$, unit marginal variance \\
innovation law & $\eta \sim \mathrm{Laplace}(0, b)$, $b = \sqrt{\ivar/2}$, excess kurtosis $+3$ \\
chain length & $\nsites = 32$ sites \\
channel & variance preserving, $x = e^{-t}a + \sqrt{\Deltat}\,\varepsilon$ \\
noise schedule & $12$ levels, geometric, $t \in [0.05,\ 3.0]$, ratio $1.45$ \\
reference score & grid BP under the true kernel, $\ngrid = 401$, $A = 8$ \\
fitted kernel & $C = 8$ mixture components, $24$ free parameters \\
training sizes & $\nseq \in \{32, \dots, 4096\}$ (headline) or $\{32,128,512,2048\}$ \\
replicates & $16$ seeds, both arms re-fitted per cell \\
validation / test & $1{,}024$ / $2{,}048$ sequences (structured runs); $256$ / $256$ (headline) \\
\bottomrule
\end{tabular}
\end{minipage}
\end{center}

\section{The estimators being compared}\label{sec:em-estimators}

\subsection{From a denoiser to a score}

Both network arms are trained by denoising score matching and both are read as
scores, so the choice of output head is a
parameterisation of the same object. If the network predicts the noise,
$\widehat{\varepsilon}_\phi$, or the clean signal, $\widehat{a}_\phi$, then
\begin{equation}
  \widehat{\score}_\varepsilon(x,t) = -\frac{\widehat{\varepsilon}_\phi(x,t)}{\sqrt{\Deltat}},
  \qquad
  \widehat{\score}_{a}(x,t) = \frac{e^{-t}\,\widehat{a}_\phi(x,t) - x}{\Deltat},
  \label{eq:em-parameterisation}
\end{equation}
and the two agree identically under the change of variable
$\widehat{a} = e^{t}(x - \sqrt{\Deltat}\,\widehat{\varepsilon})$. As
\emph{function classes} they are therefore the same; as \emph{optimisation
problems} they are not, because the substitution rescales the residual by
$e^{t}\sqrt{\Deltat}$ and so reweights the noise levels against one another.
Both are offered to every network arm, and the choice is made on validation.
The screen of Section~\ref{sec:em-screen} settles the question empirically and
without ambiguity: within each architecture the best noise-predicting
configuration beats every signal-predicting one under the pooled selection
risk ($0.148$ against $0.224$ for the window head, $0.226$ against $0.427$
for the dilated stack, $0.187$ against $0.262$ for the bidirectional
recurrence). Where a parameterisation is quoted below, it is the noise head.

Throughout, $\tau(t) = \big(t,\ e^{-t},\ \sqrt{\Deltat}\big)$ is the
three-dimensional time embedding, broadcast to every site so that a local
kernel can see it inside its own receptive field; a time feature supplied only
at the sequence ends would be invisible to a convolution.

\subsection{The structured estimator}

EM--BP fits the mixture kernel of \eqref{eq:em-mixture-kernel} by the procedure
of Section~\ref{sec:em-algorithm} and evaluates the score through the
score--posterior identity. Its free parameters are the $C-1$ free mixture
weights, $C$ component means, $C$ component variances and the autoregressive
coefficient, so at $C = 8$
\begin{equation}
  \dim\theta \;=\; (C-1) + C + C + 1 \;=\; 3C \;=\; 24,
  \label{eq:em-dim}
\end{equation}
independent of $\nsites$. Three structural facts are supplied and not learned:
the Markov factorisation, the linear-plus-innovation form of the kernel, and
homogeneity across sites. Everything else --- the innovation law itself, and
$\corr$ --- is estimated from the same noised data the networks see.

\subsection{The unstructured denoiser}

The baseline with no structural prior is a fully connected $\tanh$ network
mapping the whole sequence and the time embedding to the whole output,
\begin{equation}
  f_\phi(x,t) \;=\; W_3\,\tanh\!\big(W_2\,\tanh(W_1 [x;\tau(t)] + b_1) + b_2\big) + b_3,
  \label{eq:em-mlp}
\end{equation}
with $W_1 \in \R^{128\times(\nsites+3)}$, $W_2 \in \R^{128\times128}$,
$W_3 \in \R^{\nsites\times128}$, giving $25{,}248$ free parameters, a thousand
times \eqref{eq:em-dim}. Nothing in \eqref{eq:em-mlp} knows that the sites are
ordered: permuting the coordinates of $x$ and the rows of the output is absorbed
by permuting the columns of $W_1$ and the rows of $W_3$. The network must
rediscover the chain from data, and the comparison in
Section~\ref{sec:em-pointwise} measures what that costs.

\subsection{The three structure-aware architectures}\label{sec:em-archs}

Beating \eqref{eq:em-mlp} says little, because a large part of the margin could
be permutation-blindness rather than the Markov model. Three architectures that
do encode sequence structure are therefore screened, spanning the receptive
field from bounded to unbounded.

\paragraph{Weight-shared window head, radius $r$.} One predictor applied at
every site to a window of its neighbours,
\begin{equation}
  \widehat{y}_k \;=\; g_\phi\big(x_{k-r},\dots,x_{k+r};\ \tau(t)\big),
  \qquad k = 1,\dots,\nsites,
  \label{eq:em-window}
\end{equation}
with $g_\phi$ a two-layer $\tanh$ network of width $W$, the same $\phi$ at every
$k$, and zero padding outside the sequence. This is exactly a one-dimensional
convolution of kernel width $2r+1$ followed by pointwise layers, so it supplies
locality \emph{and} translation invariance and its parameter count,
$(2r+4)W + W^2 + 3W + 1$ --- input layer $(2r+4)W + W$, hidden layer
$W^2 + W$, scalar readout $W + 1$, matching the implementation's own count ---
does not grow with $\nsites$. Its receptive radius is
$r$ by construction: no amount of data lets it propagate information further.

\paragraph{Dilated convolutional stack.} A residual stack of three-tap
convolutions at geometrically increasing dilation,
\begin{equation}
  h^{(\ell)}_k \;=\; h^{(\ell-1)}_k
  \;+\; P_\ell \tanh\!\Big(\textstyle\sum_{s \in \{-d_\ell,\,0,\,d_\ell\}} W_{\ell,s}\, h^{(\ell-1)}_{k+s} + b_\ell\Big),
  \qquad d_\ell \in \{1,2,4,8\},
  \label{eq:em-dilated}
\end{equation}
initialised from $h^{(0)} = \tanh(W_{\mathrm{in}}[x;\tau(t)])$ and read out by a
pointwise map. Its receptive radius is $\sum_\ell d_\ell = 15$, which spans an
$\nsites = 32$ chain end to end.

\paragraph{Bidirectional message passing.} Forward and backward recurrences with
shared local updates and a joint readout,
\begin{equation}
  \overrightarrow{h}_k = \tanh\!\big(A z_k + F \overrightarrow{h}_{k-1} + a\big),
  \quad
  \overleftarrow{h}_k = \tanh\!\big(B z_k + G \overleftarrow{h}_{k+1} + b\big),
  \quad
  \widehat{y}_k = C\big[\overrightarrow{h}_k; \overleftarrow{h}_k; z_k\big] + c,
  \label{eq:em-bimp}
\end{equation}
where $z_k = [x_k; \tau(t)]$. This is the same topology as the two-sweep
sum--product schedule of Section~\ref{sec:sum-product}: one local rule applied
at every site, run once in each direction, then combined --- with affine-plus-%
$\tanh$ updates in place of the exact message updates. Of the three it has the
best claim to being able to \emph{learn} what belief propagation computes, and
its recurrent blocks are initialised orthogonally so that a perturbation at one
end of the chain is not attenuated to nothing before it reaches the other.

\begin{center}
\begin{minipage}{\linewidth}\small\centering
\captionof{table}[The ladder of structural assumptions]{The five arms as a
ladder of supplied structure, and what each is allowed to observe. ``Order''
means the architecture knows the sites are sequentially ordered; ``sharing''
means one rule is reused across sites; ``radius'' is the furthest a site's
prediction can depend on evidence. Only EM--BP is told the generating family,
and only EM--BP is denied the clean trajectory: every network arm is trained on
paired $(a,x)$ data as \eqref{eq:sm-dsm} requires, redrawing $t$ and the
corrupting noise each gradient step, while EM--BP sees one frozen noised
realisation per trajectory. The parameter ranges are those swept in
the screen of Section~\ref{sec:em-screen}; the window head at its selected
setting ($r = 2$, $W = 64$) carries $4{,}801$.}
\label{tab:em-architectures}
\begin{tabular}{@{}lcccccr@{}}
\toprule
estimator & order & sharing & kernel form & radius & observes & free parameters \\
\midrule
EM--BP, \eqref{eq:em-mixture-kernel} & yes & yes & yes & $\nsites$ & $x$ & $24$ \\
window head, \eqref{eq:em-window} & yes & yes & no & $r$ & $(a,x)$ & $4.8$k--$21.4$k \\
dilated stack, \eqref{eq:em-dilated} & yes & yes & no & $15$ & $(a,x)$ & $66.4$k--$264$k \\
message passing, \eqref{eq:em-bimp} & yes & yes & no & $\nsites$ & $(a,x)$ & $9.0$k--$34.3$k \\
global MLP, \eqref{eq:em-mlp} & no & no & no & $\nsites$ & $(a,x)$ & $25{,}248$ \\
\bottomrule
\end{tabular}
\end{minipage}
\end{center}

\noindent Table~\ref{tab:em-architectures} sets the five arms side by side as a
ladder, which is the object the rest of the chapter climbs. All network arms are
trained with Adam at batch size $64$ and gradient
clipping at $1$, over a shared learning-rate grid and a shared checkpoint ladder
from $500$ to $64{,}000$ steps, with both output heads of
\eqref{eq:em-parameterisation} offered to each. Each site contributes one
training example to a weight-shared arm, so a minibatch of $B$ sequences
supplies $B\nsites$ examples: sharing is what makes a small receptive field
cheap as well as structured.

\section{Pointwise data efficiency}\label{sec:em-pointwise}

Against the unstructured denoiser \eqref{eq:em-mlp}, and under the structural
assumptions of Table~\ref{tab:em-protocol}, the structured estimator attains
lower relative denoising error by a factor of $\ratiolo$--$\ratiohi$ across the
\nsizesused{} training-set sizes (Table~\ref{tab:pointwise},
Figure~\ref{fig:em-value-of-structure}a). It is more accurate in all but $\nreversed$ of
the $\ncellsused$ scored (seed, size, level) cells. The factor is not
monotone in the budget; Appendix~\ref{app:aggregation} reports the same
cells under four alternative summaries, of which the retired
mean-of-per-level-ratios is systematically the largest.
%% Included by the paper and the thesis (Chapter 10). The
%% workshop does not carry it: at four pages it presents the analytical result,
%% and a comparison whose whole content is a protocol cannot be defended in the
%% space available.
%%
%% GENERATED by tools/make_tab_efficiency.py from outputs/frozen/
%% exp_07_certified_seed*/ (nseq <= 2048) and exp_07_n4096_seed*/
%% (16 seeds, 1520 cells). The budget calibration in the caption comes
%% from exp_07_budget2048_seed*/. Do not hand-edit the numbers; rerun the generator.

\begin{table}[!htbp]
\caption[Relative denoising error against training-set size]{Schedule-pooled relative score error against the deterministic grid-BP
reference under the true kernel: squared error and reference score energy are summed over the
twelve noise levels before the square root, per seed; 16 seeds $\pm$ one standard error.
Both arms select their optimisation budget on a disjoint validation bundle. EM--BP is
more accurate in $1{,}518$ of $1{,}520$ (seed, size, level) cells --- a
descriptive count, not $\,1{,}520$ replications, since the twelve levels of a
cell share a training set and a fitted model and the inferential unit stays the
$16$ paired seeds. The ratio
column is the paired per-seed ratio of the two pooled errors, averaged over seeds,
with a percentile bootstrap over those ratios beside it (seeds resampled, not
cells) --- close to, but
not identical with, the quotient of the two columns beside it;
Appendix~\ref{app:aggregation} reports the same cells under four alternative aggregations.
Protocol, cell exclusions and the budget calibration are in \provenanceappendix.}
\label{tab:pointwise}
\centering
\setlength{\tabcolsep}{4pt}
\begin{tabular}{rcccc}
\toprule
sequences $\nseq$ & network & EM--BP ($24$ free) & ratio & 95\% CI \\
\midrule
32   & $0.6024 \pm 0.0046$ & $\mathbf{0.0896 \pm 0.0043}$ & $6.9 \pm 0.3$ & $[6.4,\, 7.5]$ \\
64   & $0.5025 \pm 0.0042$ & $\mathbf{0.0841 \pm 0.0066}$ & $6.6 \pm 0.5$ & $[5.6,\, 7.6]$ \\
128  & $0.4238 \pm 0.0024$ & $\mathbf{0.0572 \pm 0.0037}$ & $7.9 \pm 0.5$ & $[7.0,\, 8.8]$ \\
256  & $0.3494 \pm 0.0017$ & $\mathbf{0.0449 \pm 0.0027}$ & $8.2 \pm 0.4$ & $[7.3,\, 9.0]$ \\
512  & $0.2891 \pm 0.0010$ & $\mathbf{0.0375 \pm 0.0022}$ & $8.2 \pm 0.5$ & $[7.2,\, 9.2]$ \\
1024 & $0.2417 \pm 0.0010$ & $\mathbf{0.0294 \pm 0.0025}$ & $9.0 \pm 0.7$ & $[7.7,\, 10.4]$ \\
2048 & $0.2033 \pm 0.0007$ & $\mathbf{0.0258 \pm 0.0019}$ & $8.6 \pm 0.7$ & $[7.4,\, 9.9]$ \\
4096 & $0.1790 \pm 0.0007$ & $\mathbf{0.0180 \pm 0.0015}$ & $10.9 \pm 0.8$ & $[9.4,\, 12.5]$ \\
\bottomrule
\end{tabular}
\end{table}


\begin{figure}[t]\centering
\includegraphics[width=\linewidth]{fig_value_of_structure.pdf}
\caption[What supplying the chain is worth, against two baselines]{The two
comparisons, and the difference between them.
\emph{(a)} Headline protocol: schedule-pooled relative score error
\eqref{eq:em-pooling} against
training-set size for EM--BP and a denoiser given no structure at all, $16$
seeds, with the fitted finite-range slopes. \emph{(b)} All-site protocol: the
same quantity for EM--BP and the screened window head, which \emph{is} given
locality and weight sharing; the dotted curve is the two-parameter fit
\eqref{eq:em-floor} to the window head, whose fitted floor,
$\figwinfloor$, is indistinguishable from its measured risk at
$\nseq = 2048$. Read that as diminishing returns over the four tested
sizes and not as a resolved floor: a floorless power law
$\nseq^{-\figwinpowerslope}$ fits the same four points with a residual sum
only $\figwinrssratio$ times larger at equal parameter count, so two
residual degrees of freedom cannot separate a converged curve from a slowly
decaying one, and a fitted floor that lands on the last measured point is
partly a tautology. EM--BP's curve, by contrast, is still falling steeply.
The two panels
are not on a common scale --- they use
different scoring regions and different test bundles --- and no number is
carried between them. \emph{(c)} The two error ratios. The upper curve is the
headline $\ratiolo$--$\ratiohi\times$; the lower one is the
$\structratiolo$--$\structratiohi\times$ that survives once locality and weight
sharing are supplied, and the gap between them is how much of the headline is
not attributable specifically to the Markov model. Both grow across the tested
sizes, which
is a description of eight and four measured points respectively and not a
mechanism.}
\label{fig:em-value-of-structure}
\end{figure}

The growth of the ratio has a one-line reading that the panel supports directly.
If each arm's error follows a power law $\mathcal{E} \approx c\,\nseq^{-\varrho}$
over the tested range, then
\begin{equation}
  R(\nseq) \;=\; \frac{c_{\mathrm{base}}}{c_{\mathrm{EM}}}\,
  \nseq^{\,\varrho_{\mathrm{EM}} - \varrho_{\mathrm{base}}},
  \label{eq:em-rate-ratio}
\end{equation}
so the ratio grows precisely when the structured arm converges faster in data.
The fitted finite-range slopes of the pooled error curves are
$\varrho_{\mathrm{EM}} = \pooledslopeem \pm \pooledslopeemse$ and
$\varrho_{\mathrm{base}} = \pooledslopebase \pm \pooledslopebasese$
(log--log least squares over the
\nsizesused{} sizes --- descriptive fits, not asymptotic rates), and
\eqref{eq:em-rate-ratio} then predicts
that a $128$-fold increase in data multiplies the ratio by
$128^{\pooledslopeem - \pooledslopebase} = \pooledgrowthpred$; the measured ratio moves from
$\ratiofirst$ to $\ratiolast$, a factor of $\ratiogrowth$,
non-monotonically. Two remarks travel with that. Neither slope is the parametric
$\varrho = 1/2$, so neither arm is behaving like a correctly specified
finite-dimensional estimator over this range --- for EM--BP the mixture is
misspecified against a Laplace innovation and the grid is finite; for the
network the function class is unrestricted. And \eqref{eq:em-rate-ratio}
describes eight points; it identifies no mechanism.

Every tunable quantity either arm is allowed to choose is chosen on a validation
bundle disjoint from both training and test, and the two arms are given the same
privilege: the network selects its output head \emph{and} its training length,
EM--BP its iteration count. That choice coincides with the test-set optimum in
$\netagree\%$ of cells for the network and $\emagree\%$ for EM--BP. 

\section{How much of the margin is the architecture?}\label{sec:em-architecture}

\subsection{The bulk posterior, in closed form}\label{sec:em-locality-prediction}

Everything in this subsection is exact for the variance-matched Gaussian chain
and is used here as a \emph{design calculation} for the Laplace chain, whose
second-order structure it shares. How wide a
window does a score head need? We tried to answer with arithmetic before any
network is trained.

\paragraph{Step 1: the channel adds only to the diagonal.} The prior is
$a_k = \corr a_{k-1} + \eta_k$ with $\eta_k \sim \Nd(0,\ivar)$ and
$\ivar = 1-\corr^2$, so the marginal variance is one and
$\Sigz = (\corr^{|i-j|})$. Its precision $Q_0 = \Sigz^{-1}$ is tridiagonal ---
the Markov property written in matrix form ---
\begin{equation}
  [Q_0]_{kk} = \frac{1+\corr^2}{\ivar}\ \ \text{(interior)},
  \qquad
  [Q_0]_{k,k\pm1} = -\frac{\corr}{\ivar},
  \qquad
  [Q_0]_{kj} = 0 \ \ (|k-j| \geq 2).
  \label{eq:em-prior-precision}
\end{equation}
The channel contributes $-\tfrac{1}{2\Deltat}\norm{x - e^{-t}a}^2$ to the
log-posterior. Collecting the quadratic and linear terms in $a$,
\begin{equation}
  -\log p(a \mid x) \;=\; \tfrac12\, a^{\T} J a \;-\; h^{\T} a \;+\; \text{const},
  \qquad
  J \;=\; Q_0 + \frac{e^{-2t}}{\Deltat}\,I,
  \qquad
  h \;=\; \frac{e^{-t}}{\Deltat}\,x,
  \label{eq:em-posterior-precision}
\end{equation}
so the posterior is Gaussian with mean $m = \E[a\mid x] = J^{-1}h$. The
structural fact that drives the whole subsection is visible in
\eqref{eq:em-posterior-precision}: \emph{the channel enters the diagonal
alone.} It stiffens every site against its own observation and leaves the
nearest-neighbour couplings exactly as the prior set them. In the bulk --- away
from the two ends, where \eqref{eq:em-prior-precision} loses one neighbour ---
$J$ is therefore tridiagonal Toeplitz,
\begin{equation}
  J \;=\; \operatorname{tridiag}(\beta,\, J_{\mathrm{d}},\, \beta),
  \qquad
  J_{\mathrm{d}} \;=\; \underbrace{\frac{e^{-2t}}{\Deltat}}_{\text{evidence}}
  \;+\; \underbrace{\frac{1+\corr^2}{\ivar}}_{\text{prior}},
  \qquad
  \beta \;=\; -\frac{\corr}{\ivar}.
  \label{eq:em-bulk-params}
\end{equation}
Only $J_{\mathrm{d}}$ depends on $t$, and it is monotone: it diverges like
$1/(2t)$ as $t \downarrow 0$ and falls to the prior value as $t \to \infty$.

\paragraph{Step 2: inverting a tridiagonal Toeplitz matrix.} A matrix with three
bands has an inverse that is dense, but on the bulk of a long chain that dense
inverse is a single geometric sequence.

\begin{proposition}[Bulk covariance of the chain posterior]
\label{prop:em-bulk}
In the interior of a long chain the inverse of \eqref{eq:em-bulk-params} is
\begin{equation}
  \big(J^{-1}\big)_{k,\,k+d} \;=\; V\,q^{|d|},
  \qquad
  q \;=\; \frac{J_{\mathrm{d}} - \sqrt{J_{\mathrm{d}}^2 - 4\beta^2}}{2|\beta|} \in (0,1),
  \qquad
  V \;=\; \frac{1}{\sqrt{J_{\mathrm{d}}^2 - 4\beta^2}},
  \label{eq:em-bulk-cov}
\end{equation}
and the square root is real for every $|\corr| < 1$ and every $t > 0$.
\end{proposition}

\begin{toolbox}{Transfer-matrix inversion of a tridiagonal Toeplitz matrix}
Write $G = J^{-1}$ and look for a solution of the defining relation
$\sum_j J_{ij}G_{j,k} = \delta_{ik}$ in geometric form,
$G_{i,i+d} = V q^{|d|}$, symmetric in $\pm d$. With $i$ interior, the sum has
only three terms, and setting $k = i+d$ with $d \geq 1$ gives
\begin{equation*}
  \beta\, V q^{\,d+1} \;+\; J_{\mathrm{d}}\, V q^{\,d} \;+\; \beta\, V q^{\,d-1} \;=\; 0
  \quad\Longleftrightarrow\quad
  \beta q^2 + J_{\mathrm{d}}\, q + \beta \;=\; 0 ,
\end{equation*}
a quadratic that is the same for every $d$ --- which is why one geometric
sequence suffices. Its two roots multiply to $\beta/\beta = 1$, so they are
$q$ and $1/q$: \emph{exactly one} lies inside the unit disc, and it is the one
above. (For $\corr > 0$ all entries are positive; for $\corr < 0$ they alternate
in sign with the same modulus.) The remaining relation, at $d = 0$, fixes the
prefactor: $J_{\mathrm{d}} V + 2\beta V q = 1$, and since
$2\beta q = -\big(J_{\mathrm{d}} - \sqrt{J_{\mathrm{d}}^2 - 4\beta^2}\big)$ this
reads $V\sqrt{J_{\mathrm{d}}^2 - 4\beta^2} = 1$.

Reality of the root needs $J_{\mathrm{d}} > 2|\beta|$, and the margin is
strictly positive at every operating point,
\begin{equation*}
  J_{\mathrm{d}} - 2|\beta|
  \;=\; \frac{e^{-2t}}{\Deltat} + \frac{1 + \corr^2 - 2|\corr|}{\ivar}
  \;=\; \frac{e^{-2t}}{\Deltat} + \frac{(1 - |\corr|)^2}{\ivar}
  \;>\; 0 ,
\end{equation*}
both terms being non-negative and the first strictly positive for $t > 0$. So
$q$ is real, and $q < 1$ strictly. \hfill$\square$
\end{toolbox}

\noindent Two checks. The formula is verified against a direct numerical
inverse at $\nsites = 32$, $\corr = 0.85$: the relative disagreement over
$d = 0,\dots,5$ is $6\times10^{-15}$ at $t = 0.05$ and $2.7\times10^{-4}$ at
$t = 3$, the growth being the two boundaries leaking into the ``bulk'' once the
correlation length is an appreciable fraction of the chain --- which is the one
regime in which \eqref{eq:em-bulk-cov} should be read as approximate here. And
the limits are the ones physical sense demands,
\begin{equation}
  \begin{aligned}
  t \to 0:&\quad J_{\mathrm{d}} \to \infty,
  &&q \approx \frac{|\beta|}{J_{\mathrm{d}}} \to 0,
  &&V \to 0
  &&\text{(sites pinned by their observations)};\\
  t \to \infty:&\quad J_{\mathrm{d}} \to \frac{1+\corr^2}{\ivar},
  &&q \to |\corr|,
  &&V \to 1
  &&\text{(posterior}\to\text{prior, unit variance)},
  \end{aligned}
  \label{eq:em-q-limits}
\end{equation}
the second because $J_{\mathrm{d}}^2 - 4\beta^2 = (1-\corr^2)^2/\sigma_\eta^4 = 1$
once the evidence term vanishes. As diffusion time runs, the posterior
correlation length climbs monotonically from zero to the prior's own.

\paragraph{Step 3: the influence of distant evidence is exactly geometric.}
Because $m = \frac{e^{-t}}{\Deltat}J^{-1}x$ is linear in the observation,
Proposition~\ref{prop:em-bulk} gives the sensitivity of one site's posterior
mean to an observation $d$ sites away in closed form,
\begin{equation}
  \frac{\partial m_k}{\partial x_{k \pm d}}
  \;=\; \frac{e^{-t}}{\Deltat}\,\big(J^{-1}\big)_{k,\,k\pm d}
  \;=\; \frac{e^{-t}}{\Deltat}\, V\, q^{\,d}
  \;=\; \frac{e^{-t} V}{\Deltat}\, e^{-d/\xi},
  \qquad
  \xi \;\coloneqq\; \frac{1}{\log(1/q)} .
  \label{eq:em-influence}
\end{equation}
This is the object a score head has to represent, and it is a pure exponential
in the separation with a single length $\xi(\corr,t)$. Nothing about it is
fitted, and within the infinite-chain bulk nothing about it is asymptotic in
$t$ either; on a finite chain it is the bulk approximation, exact in the
interior and corrected near the two ends, as Chapter~\ref{ch:gaussian-bp}
derives it.

\paragraph{Step 4: what a radius-$r$ window must discard.} A head of radius $r$
cannot use $x_{k\pm d}$ for $d > r$, whatever it learns. In squared weight,
using \eqref{eq:em-influence} and two geometric sums,
\begin{equation}
  \frac{\sum_{|d| > r} q^{2|d|}}{\sum_{d \in \mathbb{Z}} q^{2|d|}}
  \;=\;
  \frac{2q^{2(r+1)}/(1-q^2)}{(1+q^2)/(1-q^2)}
  \;=\;
  \frac{2\,q^{2(r+1)}}{1+q^2} ,
  \label{eq:em-discarded}
\end{equation}
exactly. Pooling over the schedule with the closed-form weights of
Table~\ref{tab:em-locality} gives one number per radius, the
\emph{coefficient-tail fraction}
\begin{equation}
  \mathcal{T}(r) \;=\; \sum_{t \in \mathcal{T}} w_t \;\frac{2\,q(\corr,t)^{2(r+1)}}{1 + q(\corr,t)^2} .
  \label{eq:em-truncation}
\end{equation}
What \eqref{eq:em-truncation} is, and is not, has to be exact.
It is a fraction of squared coefficient mass, and it sets the geometric
\emph{exponent} of everything local. It is \emph{not} a prediction risk, for
three reasons: the omitted observations are correlated with the retained
ones, so their squared coefficients do not add to an error energy; a genuine
windowed estimator reoptimises its retained coefficients --- it is the
posterior mean of the window model, not the exact coefficients with the tail
set to zero; and a squared-mass fraction is not commensurable with a relative
RMS error even in scale. It is used below only as a decay diagnostic.

The valid risk anchor exists in closed form and costs nothing. The
Bayes-optimal radius-$r$ estimator of the design model is the windowed
posterior mean $\E[a_k \mid x_{k-r},\dots,x_{k+r}]$ --- exactly the
windowed inference whose price Chapter~\ref{ch:gaussian-bp} computes --- and
its population relative score error, pooled over the schedule with the same
weights, is
\begin{equation}
  \mathcal{W}(r)
  \;=\;
  \Bigg(
  \frac{\sum_{t} \operatorname{tr}\big[(M_r(t)-\Qt)\,\Sigt\,(M_r(t)-\Qt)^{\T}\big]}
       {\sum_{t} \operatorname{tr}\Qt}
  \Bigg)^{1/2},
  \label{eq:em-oracle}
\end{equation}
$-M_r(t)x$ being the windowed score estimator. This is the exact
local-oracle error of the Gaussian design calculation: no learning, no
sampling, the best any radius-$r$ estimator can do \emph{on the
covariance-matched Gaussian model}. Numerically, at this $(\corr,\nsites)$
and schedule,
\begin{center}
\begin{tabular}{@{}lccccc@{}}
\toprule
$r$ & $0$ & $1$ & $2$ & $4$ & $8$ \\
\midrule
$\mathcal{W}(r)$ & $0.657$ & $0.202$ & $0.089$ & $0.029$ & $0.007$ \\
\bottomrule
\end{tabular}
\end{center}
It is a design anchor, not the experiment's own floor: the experiment is
Laplace, whose local oracle is a different quantity. How different is a
question of fact rather than of ordering theory --- no result here ranks the
Bayes risks of two priors that share a covariance --- so it is measured.
The radius-$r$ Laplace oracle is
$\E[a_k \mid x_{k-r:k+r}]$ under the chain of \eqref{eq:l-chain-prior},
which is exactly the $L^2$-optimal radius-$r$ estimator of $\score_k$ because
$x_k$ lies inside the window; window-restricted grid BP over the chain's true
state marginals evaluates it, and Table~\ref{tab:em-oracle-laplace} pools it
exactly as \eqref{eq:em-oracle} pools the Gaussian one.

\begin{center}
\begin{minipage}{\linewidth}\small\centering
\captionof{table}[The two local oracles]{The exact Gaussian windowed-oracle
error \eqref{eq:em-oracle} and the Laplace chain's own, both pooled over the
schedule at $\corr = 0.85$, $\nsites = 32$. The Laplace column is Monte Carlo
over $\laporacledraws$ chains with a bootstrap standard error over draws; the
Gaussian column is closed form. The Laplace oracle is the larger of the two
wherever the difference is resolved --- at $r = 2$ by
$\laporacletwosigmas$ bootstrap standard errors --- and the two are
indistinguishable by $r = 8$.}
\label{tab:em-oracle-laplace}
\begin{tabular}{@{}lccccc@{}}
\toprule
$r$ & $0$ & $1$ & $2$ & $4$ & $8$ \\
\midrule
Gaussian $\mathcal{W}(r)$ & $0.657$ & $0.202$ & $0.089$ & $0.029$ & $0.007$ \\
Laplace $\mathcal{W}_{\mathrm{L}}(r)$
  & $\laporaclezero$ & $\laporacleone$ & $\laporacletwo$
  & $\laporaclefour$ & $\laporacleeight$ \\
\bottomrule
\end{tabular}
\end{minipage}
\end{center}

\noindent
So $\mathcal{W}(r)$ tells us what the window is worth in the second-order
geometry the two models share, and the measurement says the experiment's own
floor sits slightly above it at every radius that matters here.

\begin{center}
\begin{minipage}{\linewidth}\small\centering
\captionof{table}[Posterior correlation length across the noise schedule]{The
influence-decay rate of \eqref{eq:em-bulk-cov}, the correlation length
$\xi$ it implies through \eqref{eq:em-influence}, and the weight
\eqref{eq:em-weights} each level carries in the pooled risk, at
$\corr = 0.85$ (eight of the twelve levels shown). The two right-hand columns
pull in opposite directions, and that opposition is the whole design argument:
evidence travels furthest exactly where the pooled score error is least
sensitive. The four smallest levels carry $53.6\%$ of the weight and have
$\xi \leq 1.07$; the four largest carry $19.2\%$ and have $\xi \geq 3.55$.}
\label{tab:em-locality}
\begin{tabular}{@{}rccc@{}}
\toprule
$t$ & $q(\corr,t)$ & $\xi = 1/\log(1/q)$ & weight $w_t$ (\%) \\
\midrule
$0.0500$ & $0.203$ & $0.63$ & $16.7$ \\
$0.0725$ & $0.259$ & $0.74$ & $14.4$ \\
$0.1053$ & $0.323$ & $0.88$ & $12.2$ \\
$0.1527$ & $0.392$ & $1.07$ & $10.2$ \\
$0.3215$ & $0.541$ & $1.63$ & $7.1$ \\
$0.6769$ & $0.689$ & $2.69$ & $5.4$ \\
$1.4250$ & $0.806$ & $4.65$ & $4.8$ \\
$3.0000$ & $0.848$ & $6.06$ & $4.7$ \\
\bottomrule
\end{tabular}
\end{minipage}
\end{center}

\noindent The prediction now follows without reference to any fit. Because the
pooled risk is concentrated where $\xi \lesssim 1$, a small radius is enough
where it matters: at $t = 0.1527$ a radius of $4$ discards $q^5 = 0.9\%$ of the
influence, the oracle error \eqref{eq:em-oracle} has already fallen to
$0.09$ at $r = 2$ and $0.03$ at $r = 4$, and the coefficient tail
\eqref{eq:em-truncation} confirms that what remains beyond such radii is
contributed almost entirely by the high-noise levels, where the weight is a
few percent. Enlarging the window much past the largest correlation length
that carries appreciable weight therefore removes a quantity that is already
small in the metric being optimised, while adding parameters that must be
estimated from the same data.

\begin{quote}
\emph{Prediction.} Under this protocol the best window radius is small ---
of order the largest correlation length that carries appreciable weight, a
few sites --- and the risk should \emph{rise} again well beyond it, because
past that point the receptive field is buying a quantity smaller than the
estimation error it costs. And the residual error of the selected head should
sit well above the Gaussian design oracle $\mathcal{W}(r)$ at its radius,
because at small radii the window itself is no longer what limits a learned
head.
\end{quote}

\subsection{The screen}\label{sec:em-screen}

Three candidates were screened on $\screendevseeds$ development seeds, disjoint
from the $\structseeds$ used for the comparison: the window head
\eqref{eq:em-window}, the dilated stack \eqref{eq:em-dilated} and the
bidirectional recurrence \eqref{eq:em-bimp}. Each ran over its own
hyperparameter grid, a shared learning-rate grid, both output heads and a
shared checkpoint ladder, $\screenconfigs$ configurations in all, scored on the
\screenregion{} region. \textbf{The window head won}, at radius
$\screenwinnerradius$ and width $\screenwinnerwidth$, with the
\screenrunnerup{} network next ($\screenwinnerrisk$ against
$\screenrunneruprisk$).

\input{sections/tab-screening}

\begin{figure}[t]\centering
\includegraphics[width=\linewidth]{fig_screening.pdf}
\caption[The architecture screen, and the receptive-field question]{
\emph{(a)} Every configuration of the screen, scored on the \screenregion{}
region with errors pooled over the $\screendevseeds$ development seeds and
sizes --- the selector's own accumulation; one marker per
architecture, the outlined marker being each architecture's best. The horizontal
axis is what makes the comparison readable: the dilated stack sits an order of
magnitude to the right and does not reach the window head's risk anywhere in the
screen, so the outcome is not one arm having been starved of capacity.
\emph{(b)} Risk against window radius, at both widths, with the exact
Gaussian windowed-oracle error \eqref{eq:em-oracle} on the same axis --- the
same units, unlike a coefficient-tail fraction, which is deliberately not
drawn on a risk axis. The measured curve is minimised at $r = 2$, nearly tied
with $r = 4$, and turns upward beyond, as
Section~\ref{sec:em-locality-prediction} predicts, while the oracle keeps
falling geometrically. The vertical gap between the measured heads and the
oracle at the same radius is what learning, capacity and the Laplace/Gaussian
mismatch cost within a fixed information set.}
\label{fig:em-screening}
\end{figure}

Table~\ref{tab:screening} reports the screen in full and
Figure~\ref{fig:em-screening} shows every configuration in it, which is the more
useful object: the three architectures separate by parameter count far more
cleanly than by risk, and the two arms built to propagate further occupy the
expensive half of the axis without buying anything there.

Both halves of the prediction hold. The measured radius sweep at width $64$
is $\mathbf{0.148}$, $0.149$, $0.185$, $0.214$ at $r = 2, 4, 8, 16$ under the
pooled selection risk: minimised at $2$, in a near-tie with $4$, and rising
steadily beyond --- the behaviour of a term that has stopped
buying anything and started costing estimation error. And the selected head
sits well clear of the local oracle at its own radius: $0.148$ against the
experiment's own $\mathcal{W}_{\mathrm{L}}(2) = \laporacletwo$
(Table~\ref{tab:em-oracle-laplace}), and against the Gaussian design
anchor's $\mathcal{W}(2) = 0.089$. The first of those two differences is the
one with a reading: it is what remains to be won \emph{inside} the radius-$2$
information set --- learning, capacity and optimisation together, not
separated --- and it says the receptive field is not the binding constraint
at the selected radius. Because the measured Laplace oracle exceeds the
Gaussian one here, the larger difference from $\mathcal{W}(2)$ bounds the
first from above rather than replacing it; that ordering is a measurement at
this $(\corr, \nsites)$ and schedule, not a theorem. No finer causal split of
either gap follows from this design.

Two scoping conditions travel with the outcome. It is a model-selection fact:
under this search space, protocol and development seeds, the simplest candidate
scored best. It is not evidence that long-range propagation is useless here.

\subsection{The confirmatory comparison}

Against the screened baseline --- the radius-$\screenwinnerradius$,
width-$\screenwinnerwidth$ window head, $4{,}801$ parameters --- at the
headline protocol throughout: the same
twelve noise levels and $\structseeds$ seeds as
Table~\ref{tab:pointwise}, every site scored rather than the centre alone, the
baseline's architecture and radius fixed by the screen on separate seeds, and
each arm's checkpoint chosen on a validation bundle disjoint from test --- the
surviving margin is $\structratiolo$--$\structratiohi\times$.

\input{sections/tab-structured}

Table~\ref{tab:structured} reports it, and
Figure~\ref{fig:em-value-of-structure}b and c place it beside the headline;
the dotted curve in panel (b) is the two-parameter floor fit
\begin{equation}
  \mathcal{E}(\nseq)^2 \;=\; \mathcal{E}_{\infty}^2 + \frac{c}{\nseq},
  \label{eq:em-floor}
\end{equation}
a floor plus a $1/\nseq$ estimation term, fitted to the window head's four
squared risks --- descriptive, as a four-point fit must be.
EM--BP is more accurate in every one of the $\structcells$ scored cells. That
count is $\structseeds$ seeds $\times$ $\structsizes$ sizes $\times$
$\structregions$ scoring regions (\structregionlist{}), of which the table
displays two; it is a descriptive diagnostic and not a replication count, since
regions and sizes share fitted models and training data, so the inferential unit
behind each displayed row is the $\structseeds$ paired seeds.
Appendix~\ref{app:aggregation} reports the same ratio under $\structaggcount$
aggregations; none of them, at any size, falls below $\structaggfloor\times$.

A large fraction of the headline margin therefore does not survive supplying
locality and weight sharing. That comparison does not decompose causally:
architecture, parameter sharing, capacity and optimisation all change when the
baseline is swapped, and no design here separates them. What it does establish
is that the generic-MLP ratio overstates what can be attributed specifically to
the supplied Markov model, and by a factor that runs from about three at the
smallest training set to about $1.4$ at the largest.

\section{How much mixture capacity the estimator needs}\label{sec:em-capacity}

The estimator's one nonparametric dial is the number of mixture components.
The capacity sweep refits $C \in \{1, 2, 4, 8, 16\}$ on identical data,
paired within seed against the single-Gaussian innovation $C = 1$, at
$\nseq \in \{\capsizes\}$, and scores held-out log-evidence.

\input{sections/tab-capacity}

\begin{figure}[t]\centering
\includegraphics[width=\linewidth]{fig_capacity.pdf}
\caption[Mixture capacity against a single Gaussian innovation]{\emph{(a)}
Held-out log-evidence per edge, paired within seed against $C = 1$: no tested
capacity improves on a single Gaussian innovation at either size, and at
$\nseq = \capheadsize$ the $C = 16$ deficit against $C = 8$ is resolved
against zero.
\emph{(b)} Why (a) is not a clean capacity test: the share of fits stopped at
the iteration cap, and with a narrowest component under the grid's resolution
floor, both climb steeply with $C$.}
\label{fig:em-capacity}
\end{figure}

Table~\ref{tab:capacity} and Figure~\ref{fig:em-capacity} give the outcome
and its caveat in one view. Under this protocol no capacity buys anything:
every paired mean is negative, the $\nseq = \capheadsize$ deficit of the
largest capacity against $C = 8$ is $\capheadverdict$ against zero
($\capheadlogev$, CI $\capheadlogevci$), and only two to four seeds in
sixteen ever prefer a mixture. But the sweep is not a clean statement about
capacity in the abstract: from $C = \capcappedfrom$ upward every cell stops at the
$\capiterlimit$-iteration cap rather than at tolerance, and the share of fits
whose narrowest component falls below the grid's resolution floor climbs with
$C$, so a statistical disadvantage of large mixtures is entangled with an
optimisation and a discretisation effect. The defensible conclusion is
protocol-scoped: under this grid, M-step, iteration cap and validation rule,
capacity beyond a single Gaussian innovation bought nothing --- and the
convergence economics of Chapter~\ref{ch:em-kernel} say why more capacity is
also \emph{slower}: every added component deepens the missing-information
problem the channel already creates.

That finding has a consequence for how the headline of
Section~\ref{sec:em-pointwise} may be read, and it is worth stating against
interest. The headline estimator runs at $C = 8$ because the protocol fixed it
there in advance, not because any selection rule chose it; in this sweep the
single-Gaussian innovation is at least as good on \emph{both} criteria, the
held-out evidence and the score risk the thesis actually reports. A single
Gaussian innovation with a fitted $\corr$ is, by
Theorem~\ref{thm:l-lmmse}, exactly the covariance-matched linear estimator, so
what that comparison says is that within the range tested here the margin over
the network arms is earned by the Markov factorisation and its covariance, and
not by the mixture's capacity to represent a non-Gaussian innovation shape.
The two experiments run different bundles and scoring regions, so the risks
are not comparable across them and no headline number is restated here; the
transferable point is the attribution, not a substituted figure.

\section{Where the Markov assumption stops paying}\label{sec:em-nonmarkov}

Every result so far assumes the clean chain is exactly first-order Markov. Two
controlled violations are measured against an exact reference, so that the
estimator being wrong is not confused with the reference being wrong.

\subsection{Two violations, defined}

\paragraph{A shared global latent, strength $\beta$.} One standard normal
variable is drawn per sequence and added to every site,
\begin{equation}
  a \;=\; \frac{y + \beta g}{\sqrt{1+\beta^2}},
  \qquad y \sim \text{chain prior},
  \quad g \sim \Nd(0,1),
  \label{eq:em-global}
\end{equation}
variance preserving for every $\beta$ and Markov only at $\beta = 0$. Its
covariance is
$\Sig_\beta = (\Sigz + \beta^2 \mathbf{1}\mathbf{1}^\T)/(1+\beta^2)$ with
$\Sigz$ the AR(1) covariance, so a Laplace and a Gaussian version differ
\emph{only beyond second moments} --- the controlled comparison this thesis uses
throughout. At $\beta = 1$ half the marginal variance of every site is one
shared constant.

\paragraph{Long-range precision coupling, strength $\gamma$.} Weak couplings are
added directly to the precision at every separation of two or more,
\begin{equation}
  Q_\gamma \;=\; Q_{\mathrm{AR}} + \gamma\,Q_{\mathrm{long}},
  \qquad
  [Q_{\mathrm{long}}]_{ij} = -e^{-|i-j|/\ell}\ \ (|i-j| \geq 2),
  \quad
  [Q_{\mathrm{long}}]_{ii} = 1.05 \sum_{|i-j| \geq 2} e^{-|i-j|/\ell},
  \label{eq:em-longrange}
\end{equation}
with $\ell = 8$ and the diagonal compensation making $Q_\gamma$ strictly
diagonally dominant, hence positive definite; the covariance is rescaled to unit
diagonal so the prior stays variance preserving. Not Markov for
$\gamma \neq 0$.

\subsection{Why one is partly absorbable and the other is not}

The two mechanisms are not two amounts of the same thing.

\begin{proposition}[The global latent perturbs the chain in one direction]
\label{prop:em-rankone}
Let $Q_{\mathrm{AR}} = \Sigz^{-1}$ be the AR(1) precision and
$Q_\beta = \Sig_\beta^{-1}$. Then
\begin{equation}
  Q_\beta \;=\; (1+\beta^2)\,Q_{\mathrm{AR}} \;-\; c_\beta\, u u^{\T},
  \qquad
  u = Q_{\mathrm{AR}}\mathbf{1},
  \qquad
  c_\beta = \frac{(1+\beta^2)\,\beta^2}{1 + \beta^2\,\mathbf{1}^{\T}Q_{\mathrm{AR}}\mathbf{1}},
  \label{eq:em-woodbury}
\end{equation}
so the departure from a rescaled chain precision has rank exactly one, in the
single direction $u$. Moreover the lag-one correlation of \eqref{eq:em-global},
and hence the correlation of its best first-order chain approximation, is
available in closed form:
\begin{equation}
  \corr_{\mathrm{CL}}(\beta) \;=\; \frac{\corr + \beta^2}{1+\beta^2}.
  \label{eq:em-chowliu}
\end{equation}
\end{proposition}

\begin{proof}
Sherman--Morrison applied to
$\Sig_\beta = (\Sigz + \beta^2\mathbf{1}\mathbf{1}^\T)/(1+\beta^2)$ gives
\eqref{eq:em-woodbury} directly, the rank-one term being the correction that
inverting a rank-one update produces. For \eqref{eq:em-chowliu}, the covariance
of \eqref{eq:em-global} at separation $d \geq 1$ is
$(\corr^{d} + \beta^2)/(1+\beta^2)$ and the marginal variance is $1$; at
$d = 1$ this is the stated value. The KL projection of a Gaussian onto a fixed
tree matches that tree's node and edge marginals
\citep{chow1968approximating}, so on a path graph the fitted chain reproduces
the true adjacent correlations. \hfill$\square$
\end{proof}

\noindent Numerically, \eqref{eq:em-woodbury} is exact to machine precision: the
correction has one nonzero eigenvalue at every $\beta$ tested, and its
eigenvector agrees with $u/\norm{u}$ to $|\langle\cdot,\cdot\rangle| = 1.000000$.
The direction $u$ is, up to boundary terms, the constant vector --- the chain's
mean level. By contrast the correction in \eqref{eq:em-longrange} is
$\gamma Q_{\mathrm{long}}$, dense by construction, with numerical rank
$\nsites = 32$ and all $32$ eigenvalues above $1\%$ of its largest at the
strengths tested. \emph{One direction against thirty} is the structural
difference, and a chain that may reparameterise a single scalar can trade
against the first in a way it cannot against the second --- partly: matching
the lag-one correlation absorbs the covariance change along that one
direction, but a homogeneous first-order chain cannot represent the dense
rank-one precision correction exactly, so ``absorbable'' means \emph{partly
absorbed through the refitted local parameter}, not reproduced.

The fitted coefficient does exactly that. Under \eqref{eq:em-global} the fitted
$\widehat{\corr}$ climbs from $\nmrhoclean$ on the clean chain to
$\nmrhobetamax$ at $\beta = \nmbetamax$, tracking \eqref{eq:em-chowliu}
($0.9250$ at $\beta = 1$) to within $\nmcldevbeta$. Under
\eqref{eq:em-longrange} the same gap is $\nmcldevfactor$ times larger. That is
one scalar agreeing, not the whole fitted law certified --- but it is the
mechanism the proposition predicts, appearing in the fit.

\subsection{A covariance scalar orders the tested conditions}

Proposition~\ref{prop:em-rankone} explains the qualitative difference. A single
number makes it quantitative. The scalar used is the per-site KL divergence
between the \emph{covariance-matched centred Gaussians} of the contaminated
law and of its best first-order chain,
\begin{equation}
  \delta \;=\; \frac{1}{\nsites}\,
  \kl{\Nd(0,\Sig)}{\Nd(0,\Sig_{\mathrm{CL}})}
  \;=\; \frac{1}{2\nsites}\Big[
  \operatorname{tr}\big(\Sig_{\mathrm{CL}}^{-1}\Sig\big) - \nsites
  + \log\frac{\det \Sig_{\mathrm{CL}}}{\det \Sig}\Big],
  \label{eq:em-delta}
\end{equation}
closed-form for both mechanisms and involving no fit and no data. Its status
differs by arm, and the distinction is not cosmetic. For the \emph{Gaussian}
contaminated laws, \eqref{eq:em-delta} is the exact distributional KL to the
best Gaussian chain --- the excess code length that the Gaussian Chow--Liu
projection minimises. For the \emph{Laplace} arms the same expression is only
a covariance proxy: the true laws are non-Gaussian, their KL involves entropy
terms that covariance does not determine, and two non-Gaussian laws with
identical covariances can sit at different true divergences. Where the
distinction matters we write $\delta_{\mathrm{cov}}$ for the proxy reading;
it is a second-order misspecification measure, blind to everything beyond
covariance, and that is exactly the information the chain estimator's Gaussian
design geometry uses.

\begin{center}
\begin{minipage}{\linewidth}\small\centering
\captionof{table}[Misspecification, measured in nats, orders both
mechanisms]{The advantage against the closed-form covariance distance
\eqref{eq:em-delta}
from the contaminated law to its best first-order chain --- exact KL for the
Gaussian rows, covariance proxy $\delta_{\mathrm{cov}}$ for the Laplace
column. Read against their own
strength parameters the two mechanisms look unrelated; read against $\delta$
they interleave and the sequence is monotone. The entry that carries the
argument is $\beta = 1.00$ against $\gamma = 0.05$: a prior in which half of
every site's variance is a shared global constant is \emph{closer} to the chain
family, by a factor of $1.7$ in nats, than one with the weakest long-range
coupling tested --- and the measured advantage is ordered the same way. Ratios
here are means over the twelve noise levels of the per-level error ratio (a
mean of ratios, unlike the pooled headline; the estimand is stated because the
two are not interchangeable), against the convolutional and fully connected
baselines.}
\label{tab:em-misspecification}
\begin{tabular}{@{}lcccc@{}}
\toprule
contamination & $\delta$ (nats/site) & vs CNN & vs MLP & Laplace, vs CNN \\
\midrule
none                      & $0$                   & $16.4$ & $29.3$ & $11.1$ \\
global, $\beta = 0.10$    & $4.6\times10^{-6}$    & $14.8$ & $22.5$ & $10.9$ \\
global, $\beta = 0.25$    & $1.5\times10^{-4}$    & $6.7$  & $11.3$ & $6.2$  \\
global, $\beta = 0.50$    & $1.4\times10^{-3}$    & $3.2$  & $4.0$  & $3.0$  \\
global, $\beta = 1.00$    & $6.5\times10^{-3}$    & $2.2$  & $1.7$  & $2.1$  \\
long-range, $\gamma = 0.05$ & $1.1\times10^{-2}$  & $1.2$  & $1.7$  & --- \\
long-range, $\gamma = 0.10$ & $2.2\times10^{-2}$  & $\mathbf{0.97}$ & $1.2$ & --- \\
long-range, $\gamma = 0.20$ & $4.0\times10^{-2}$  & $\mathbf{0.88}$ & $\mathbf{0.84}$ & --- \\
long-range, $\gamma = 0.40$ & $6.0\times10^{-2}$  & $\mathbf{0.75}$ & $\mathbf{0.70}$ & --- \\
\bottomrule
\end{tabular}
\end{minipage}
\end{center}

\input{sections/tab-nonmarkov}

\begin{figure}[t]\centering
\includegraphics[width=\linewidth]{fig_nonmarkov.pdf}
\caption[Where each Markov violation stops paying, and the one scalar that
orders them]{\emph{(a)} The closed-form distance \eqref{eq:em-delta} to the
chain family. The two mechanisms reach the shaded break-even band at strengths
that differ by more than an order of magnitude, which is the whole of the
difference between them. \emph{(b)} The measured advantage against strength.
Under a global latent it decays smoothly and does not reach the line even where
half the marginal variance is a shared constant; under long-range coupling it
crosses within the first tested strength. The Laplace and Gaussian curves track
each other across the whole rank-one sweep, so the result is not confined to the
Gaussian arm; the long-range panel is Gaussian only, since the contaminated
Laplace model has no closed-form reference. \emph{(c)} The same nine points
against $\delta$: the two mechanisms interleave and fall on one monotone curve,
crossing break-even between the adjacent tested strengths bracketed in \eqref{eq:em-breakeven}. All curves are
schedule means, so they show the direction rather than the worst case, and all
fall because the estimator degrades rather than because the baselines improve
--- Table~\ref{tab:nonmarkov} carries EM--BP's own error alongside.}
\label{fig:em-nonmarkov}
\end{figure}

Table~\ref{tab:em-misspecification} and Figure~\ref{fig:em-nonmarkov} report
$\delta$ beside the measured advantage. Two readings follow, and the second is
the one this chapter contributes.

Read against strength, the mechanisms behave differently in kind. At
$\beta = \nmbetamax$ the Gaussian arm is still ahead at every noise level
against both baselines, worst ratio $\nmbetaminratio$; the Laplace arm is still
ahead against the convolutional baseline, worst ratio
$\nmlapbetaminratiocnn$, and reaches break-even against the fully connected one
at the largest diffusion time, at $\nmlapbetaminratio$. Under long-range
coupling the advantage against the convolutional baseline is gone by
$\gamma = \nmcrossover$, having survived $\gamma = \nmlastholding$, and by
$\gamma = 0.2$ the estimator loses to both baselines and not to one; at
$\gamma = \nmgammamax$ its own relative score error reaches $\nmemworsterr\%$.

Read against $\delta$, the tested conditions collapse onto one empirical
curve. The nine contaminated
cells interleave --- $\beta = 1$ sits between $\gamma = 0$ and
$\gamma = 0.05$ --- and the advantage is monotone in $\delta$ across both
mechanisms, over
four orders of magnitude in $\delta$ and a factor of twenty in the ratio,
against both baselines and for both innovation laws. The crossing is reported
as the bracket the design actually resolves, between adjacent tested
strengths:
\begin{equation}
  \delta^{\star} \in \big[1.1, 2.2\big]\times 10^{-2}\ \text{nats/site (vs CNN)},
  \qquad
  \big[2.2, 4.0\big]\times 10^{-2}\ \text{nats/site (vs MLP)} ,
  \label{eq:em-breakeven}
\end{equation}
the CNN ratio falling from $1.2$ to $0.97$ across its bracket and the MLP one
from $1.2$ to $0.84$ across its own. A straight line through those endpoints
crosses at $2.1$ and $3.2 \times 10^{-2}$ respectively, and those
interpolations are quoted nowhere as results: each condition here is a single
run scored across the schedule, with no seed replication, so no interval can
be attached to a crossing and a point estimate would imply a precision the
design does not have. Three further scope conditions travel with the bracket.
The two mechanisms occupy different parts of the $\delta$ axis, so the
crossing is bracketed by different conditions in each case; the ratios in this
section are schedule means of per-level ratios rather than the pooled estimand
of \eqref{eq:em-pooling}, which changes their numerical scale though not their
ordering; and adjacent strengths differ by roughly a factor of two, which is
the resolution of the whole statement.
This collapse is an empirical regularity of the two violation paths
constructed here, not a sufficiency theorem, and there is a structural reason
to expect it to have limits: even between Gaussians, equal KL does not imply
equal score geometry --- the divergence is a spectral aggregate, while the
score error of a chain approximation is
$\operatorname{tr}[(Q_{\mathrm{chain}}-Q_{\mathrm{true}})\,\Sig\,
(Q_{\mathrm{chain}}-Q_{\mathrm{true}})^{\T}]$, which depends on \emph{where}
the precision error points. Two perturbations of matched $\delta$ whose
discrepancies live in different spatial modes could in principle separate, and
nothing here tests that. Equation \eqref{eq:em-breakeven} is interpolation
within the tested conditions, not a universal tolerance for non-Markov
structure.

\section{Why the shape coordinate is the expensive one}
\label{sec:em-shape-information}

Chapter~\ref{ch:em-kernel} found the innovation shape converging far slower than
the autoregressive coefficient, and attributed it to the fourth cumulant being
damped by the channel while the second is preserved. That is a claim about the
noised likelihood, and it can be checked directly against Fisher information,
independently of any optimiser.

\subsection{What the cumulant law predicts}

Cumulants add under independent sums and a Gaussian contributes only to the
second, so for $X_t = e^{-t}A + \sqrt{\Deltat}\,\varepsilon$,
\begin{equation}
  \kappa_2(X_t) = e^{-2t}\kappa_2(A) + \Deltat = 1,
  \qquad
  \kappa_4(X_t) = e^{-4t}\kappa_4(A),
  \qquad
  \kappa_{11}(X_{t,i}, X_{t,i+1}) = e^{-2t}\corr,
  \label{eq:em-cumulants}
\end{equation}
 Now consider what identifies each
coordinate. The autoregressive coefficient is second order: it enters through
the lag-one cumulant, which the channel attenuates by $e^{-2t}$. A shape
parameter --- one that moves the fourth cumulant at exactly fixed variance --- is
fourth order, and the channel attenuates its signature by $e^{-4t}$. Reading
information as squared sensitivity over sampling variance --- a delta-method
heuristic, not a likelihood calculation: it treats the cumulant estimator as
the inference channel and its variance as approximately constant in $t$ ---
\begin{equation}
  I_\theta \;\asymp\;
  \frac{\big(\partial_\theta \kappa(X_t)\big)^2}{\operatorname{Var}\widehat{\kappa}},
\end{equation}
and noting that the denominators are $t$-independent \emph{to leading order}
because
$\kappa_2(X_t) \equiv 1$ pins the scale of the noised variable at every $t$
(higher moments of the noised law do drift with $t$, which is part of what the
heuristic omits),
\begin{equation}
  \boxed{\;
  I_{\corr}(t) \;\asymp\; e^{-4t},
  \qquad
  I_{p}(t) \;\asymp\; e^{-8t},
  \qquad
  \frac{I_{p}(t)}{I_{\corr}(t)} \;\asymp\; e^{-4t}.
  \;}
  \label{eq:em-info-prediction}
\end{equation}
The third statement is the robust one: it is a \emph{contrast}, so whatever
common factors the leading-order reading omits cancel from it.

\subsection{The measurement}

The prediction is tested on a family in which shape moves at exactly fixed
variance. The innovation law is a generalised Gaussian,
$f_p(\eta) \propto \exp(-|\eta/s|^{p})$, with $s$ fixed by
$\operatorname{Var}(\eta) = \ivar = 1-\corr^2$ at every $p$, so $p$ is a genuine
shape coordinate and not a reparameterised scale; $p = 1$ is the Laplace chain
of Chapter~\ref{ch:laplace}, $p = 2$ Gaussian. For
$\theta = (\corr, \ivar, p)$, Fisher's identity \eqref{eq:em-fisher} supplies
the exact per-sequence score of the marginal log-likelihood from one
forward--backward pass, and the information matrix is estimated in
outer-product form $J = D^{-1}\sum_d g_d g_d^{\T}$ over $D = 512$ sequences of
$\nsites = 32$ sites at $\ngrid = 401$, $A = 8$. The quantity reported is the
\emph{efficient} information $I_{p\cdot(\corr,\ivar)} = 1/(J^{-1})_{pp}$, what
remains about $p$ once the nuisance coordinates must be estimated too --- equal
to the diagonal entry only when the coordinates are orthogonal and strictly
smaller otherwise.

Two things follow from that construction and travel with every number below.
An outer-product matrix is the Fisher information only at the true parameter
of a correctly specified model, and it is here a $D$-sample estimate of it, so
these are estimates with sampling error and not an exact information
calculation. And the reported quantity is an \emph{inverse}, so its meaning
degrades with the conditioning of $J$ --- which the run records per cell.
That conditioning is benign exactly where the information is large and fails
exactly where it vanishes: $\kappa(J) = \infcondbest$ at
$t = \infcondbestt$ on the Laplace chain, a median of $\infcondmedian$ across
the sweep, and $\infcondworst$ in the worst cell
($t = \infcondworstt$, $p = \infcondworstp$), where the efficient information
is itself down at $\infcondworstinfo$. The decay these curves measure is
therefore resolved at the end where the coordinate is identified, and the far
tail is a statement that the information has gone, not a measurement of how
much remains.

\begin{figure}[t]\centering
\includegraphics[width=\linewidth]{fig_information.pdf}
\caption[Fisher information against the cumulant prediction]{
\emph{(a)} Efficient information for each coordinate against diffusion time, at
Laplace innovations, with the predicted rates \eqref{eq:em-info-prediction} as
dotted guides. On a linear time axis an exponential is a straight line, so the
comparison is read from the slopes. Both measured decays are shallower than
predicted by a common offset, which the leading-order reading does not capture.
\emph{(b)} The contrast, which is the prediction that survives the offset: the
ratio $I_p/I_{\corr}$ for all five shape exponents, normalised at
$t = 0.05$, against $e^{-4t}$. The five curves bracket the prediction, with a
mean fitted log-slope of $-3.72$ against $-4$.}
\label{fig:em-information}
\end{figure}

\begin{center}
\begin{minipage}{\linewidth}\small\centering
\captionof{table}[Fisher information lost to the channel, by coordinate]{How
much efficient Fisher information about each coordinate the channel destroys
between the smallest and largest diffusion time tested ($t = 0.05$ against
$t = 1.6$), and the contrast \eqref{eq:em-info-prediction} predicts. Entries are
ratios of outer-product estimates at $D = 512$ sequences, so only the order is
meaningful. The last
column is the quantity the prediction addresses: its geometric mean over the
five shapes is $465$ against a predicted $e^{4\times1.55} = 493$.}
\label{tab:em-information}
\begin{tabular}{@{}rcrrrr@{}}
\toprule
true $p$ & excess kurtosis & loss, $\corr$ & loss, $\ivar$ & loss, $p$ & contrast $\ \text{loss}_p/\text{loss}_{\corr}$ \\
\midrule
$1.0$ & $+3.00$ & $78\times$ & $235\times$ & $8.7\times10^{3}$ & $112$ \\
$1.2$ & $+1.74$ & $87\times$ & $392\times$ & $2.0\times10^{4}$ & $236$ \\
$1.5$ & $+0.76$ & $83\times$ & $459\times$ & $6.4\times10^{4}$ & $776$ \\
$2.0$ & $\phantom{+}0.00$ & $75\times$ & $501\times$ & $6.5\times10^{4}$ & $867$ \\
$3.0$ & $-0.58$ & $56\times$ & $522\times$ & $6.8\times10^{4}$ & $1222$ \\
\midrule
\multicolumn{5}{@{}l}{\emph{geometric mean of the contrast}} & $465$ \\
\multicolumn{5}{@{}l}{\emph{predicted by \eqref{eq:em-info-prediction}}} & $493$ \\
\bottomrule
\end{tabular}
\end{minipage}
\end{center}

Figure~\ref{fig:em-information} plots the decay against the predicted rates and
Table~\ref{tab:em-information} reads the endpoints. 

\section{What the chapter establishes}\label{sec:em-summary}

\begin{center}
\begin{minipage}{\linewidth}\small\centering
\captionof{table}[The four questions, at the strength the evidence
supports]{Each question, the finding, and the conditions the finding is
conditional on. Nothing in the right-hand column is a hedge added after the
fact; each is a stated feature of the design that a different design would have
to remove.}
\label{tab:em-verdicts}
\begin{tabular}{@{}p{0.235\linewidth}p{0.35\linewidth}p{0.33\linewidth}@{}}
\toprule
question & finding & conditional on \\
\midrule
Value of the chain, against no structure
& $\ratiolo$--$\ratiohi\times$ lower pooled error over \nsizesused{} sizes; more accurate in all but $\nreversed$ of $\ncellsused$ cells
& the schedule-pooled paired estimand of \eqref{eq:em-pooling}; both arms budget-selected on validation \\[0.4em]
Value of the \emph{model}, against locality
& $\structratiolo$--$\structratiohi\times$; EM--BP ahead in every scored cell; the baseline's error curve shows diminishing returns over four sizes (a fitted floor and a floorless power law are not separable there) while EM--BP's is still falling
& one screened architecture; four sizes; budget grids bind for both arms; no causal split of the gap \\[0.4em]
Mixture capacity
& no tested capacity beats a single Gaussian innovation; $C{=}16$ resolved \emph{worse} than $C{=}8$ at $\nseq = \capheadsize$
& this grid, M-step, iteration cap and validation rule; not a statement about capacity in the abstract \\[0.4em]
Tolerance to non-Markov structure
& advantage is monotone in one closed-form covariance scalar across the tested conditions; break-even bracketed between adjacent tested strengths, $\delta^{\star} \in [1.1,2.2]\times10^{-2}$ (CNN) and $[2.2,4.0]\times10^{-2}$ (MLP) nats/site
& nine cells, two mechanisms; exact KL only for the Gaussian arms; a measured regularity, not a derived law \\
\bottomrule
\end{tabular}
\end{minipage}
\end{center}

\noindent Table~\ref{tab:em-verdicts} states the four findings at the strength
their evidence supports. Three consequences carry into
Chapter~\ref{ch:conclusions}. Supplying the
generating family is worth a large factor at every data budget tested, but
most of the headline factor does not survive a baseline that is merely given
locality and weight sharing --- which bounds how much of it can be attributed
specifically to the Markov model, without decomposing causally what the swap
of baselines changes. Across the tested violations, the measured advantage
ordered monotonically with a covariance scalar computable before any fitting
--- an empirical regularity of the two paths probed, useful as a screening
variable and not established as a decision rule. And the coordinate on which
every non-Gaussian claim in this
thesis depends is simultaneously the slowest to converge and the poorest in
information per observation, which is why the budget disclosures in this chapter
are part of its results rather than an appendix to them.
